\section{InstructGPT 与早期 ChatGPT：效果、成本和边界}

RLHF 是否值得，不能只看损失函数，还要看用户更偏好什么、投入多少资源、原有能力是否回退。本章依次阅读三张课堂证据：偏好曲线、训练成本和 2023 年初 ChatGPT 的经验总结。

\subsection{100 倍更小，为什么仍更受偏好}

InstructGPT 论文把不同规模的 base GPT-3、SFT、PPO 与 PPO-ptx 模型放到相同的人类偏好评测中。横轴是模型参数规模，纵轴是相对某个基线被选中的频率。最醒目的结果是：1.3B 的 InstructGPT 变体在该 customer-prompt 分布上可被偏好于 175B 的 base GPT-3。

\lecturefigure{slide-08-instructgpt-preferences.jpg}{InstructGPT 偏好曲线：alignment fine-tuning 可以显著改变用户感受到的效用。}{Stanford Online 官方视频 00:09:24--00:15:18。}

\begin{knowledgebox}{读图：先看 y 轴，再看模型大小}
y 轴不是通用 intelligence，也不是所有 benchmark 的平均分，而是人类在特定 prompt 分布上的 preference win rate。图支持 ``较小 aligned model 可以更符合用户意图''；它不支持 ``1.3B 模型在知识、推理和所有任务上都胜过 175B 模型''。PPO-ptx 曲线还提醒读者，偏好收益与能力保留需要同时测量。
\end{knowledgebox}

若把偏好实验写成二项随机变量，模型 $A$ 对模型 $B$ 的估计胜率为

\[
\widehat{p}_{A>B}=\frac{1}{N}\sum_{i=1}^{N}\mathbb{1}[A\text{ 在第 }i\text{ 个比较中被选中}],
\]

其中，$N$ 是有效比较数，$\mathbb{1}[\cdot]$ 是指示函数。这个统计量依赖 prompt 分布、labeler、展示顺序和评价规范；换一批任务或人群，$\widehat{p}_{A>B}$ 可能改变。

\begin{warningbox}{Preference 不是 capability 的同义词}
偏好可以奖励清晰、合作、拒绝不当请求和遵循格式，也可能奖励自信、讨好或熟悉风格。因而 preference win 应与真实性、校准、鲁棒性、专业任务正确率和安全评测并列，而不能单独承担 ``模型更好'' 的全部含义。
\end{warningbox}

\subsection{Alignment fine-tuning 的资源账}

成本 slide 将 GPT-3 pretraining、InstructGPT SFT 与 InstructGPT RL 的 compute 放在同一纵轴，并在右侧单列人类反馈。它展示的不是 ``alignment 免费''，而是不同资源维度不对称：相对预训练，SFT/RL 的计算量很小；但收集高质量 demonstrations 与 comparisons 需要大量人类时间、规范设计和质量控制。

\lecturefigure{slide-09-instructgpt-training-costs.jpg}{InstructGPT 训练成本：alignment compute 很小，人类反馈仍约需 20,000 小时。}{Stanford Online 官方视频 00:15:18--00:19:05。}

\begin{knowledgebox}{读图：不要把 PFLOP/s-days 与美元直接相加}
左侧柱状图比较计算量，右侧文字给出讲者估算的人类反馈约 50 万美元、约 2 万小时。二者是不同资源：compute 可通过硬件并行，人类判断受到招募、培训、一致性和代表性约束。真正的系统成本还包括数据治理、评测设计、工程迭代与失败实验，slide 没有声称覆盖全部生命周期成本。
\end{knowledgebox}

\begin{table}[H]
\centering
\caption{三类资源不能只压成一个价格数字}
\begin{tabularx}{\textwidth}{L{2.9cm} X X}
\toprule
资源 & 主要作用 & 主要风险 \\
\midrule
Pretraining compute & 获得广泛语言与任务能力 & 成本高、数据边界不透明、知识截止 \\
Alignment compute & 让已有能力更可控、更符合交互目标 & reward hacking、能力回退、过拟合偏好 \\
Human feedback & 提供 demonstrations、comparisons 与错误判断 & 代表性不足、规范漂移、疲劳与一致性问题 \\
\bottomrule
\end{tabularx}
\end{table}

\subsection{ChatGPT：对话界面提高可用性，却没有消除失配}

讲者把 ChatGPT 看作 InstructGPT 的下一步：对话成为通用接口，用户可以追问、纠正和请求重写；拒绝有害任务也有所改善。但这不是 ``对齐完成''。2023 年初最突出的局限仍是 hallucination 与 prompt sensitivity：模型会编造事实，而且不同措辞可能诱发不同质量或不同安全行为。

\lecturefigure{slide-10-chatgpt-lessons.jpg}{2023 年初对 ChatGPT 的课堂总结：dialog、复杂任务、hallucination 与 prompt sensitivity。}{Stanford Online 官方视频 00:19:05--00:20:19。}

\begin{warningbox}{时间边界：不要用后来的 ChatGPT 解释这张 slide}
这页描述 2023-01-17 的公开系统。讲者在 Q\&A 明确说，ChatGPT 的精确训练数据量和步骤并未公开，只能说流程大体类似 InstructGPT；他回忆的 comparisons、demonstrations 与 20,000 小时也是近似量。讲义不能把后来的产品能力、工具调用或政策变化回填为当时已知规格。
\end{warningbox}

\begin{importantbox}{Prompt sensitivity 为什么也是 alignment 问题}
如果两个语义等价的 prompt 仅因表述方式不同就触发完全不同的真实性、拒绝或推理质量，模型学到的仍是表面模式，而不是稳定任务意图。理想系统应在保留合理个性化的同时，对语义不变换保持关键行为不变，并在歧义处主动询问。
\end{importantbox}

\subsection{本章小结}

InstructGPT 证明 alignment fine-tuning 能显著改善特定分布上的人类偏好，而且计算成本远低于预训练；但反馈劳动、代表性和 proxy 风险仍然昂贵。早期 ChatGPT 进一步证明对话接口本身能提升可用性，同时也暴露 hallucination 与 prompt sensitivity 尚未解决。

